\section*{Hoofdstuk \ref{ch:g7:moon-phases} --- Schijngestalten van de Maan}

\begin{solution}{exo:g7:moon-phases:1}
Precies de helft, altijd --- de helft aan de zonkant. Een
\omterm{def:g7:moon-phases:phases}{schijngestalte} is alleen het uitzicht: het deel van die verlichte
helft dat toevallig naar de Aarde gekeerd staat.
\end{solution}

\begin{solution}{exo:g7:moon-phases:2}
Nieuw, wassende sikkel, eerste kwartier, \omterm{def:g7:moon-phases:phases}{wassend} gibbeus,
vol, \omterm{def:g7:moon-phases:phases}{afnemend} gibbeus, laatste kwartier, afnemende
sikkel --- ongeveer $29.5$ dagen rond. \omterm{def:g7:moon-phases:phases}{Wassend} is de groeiende
helft, van nieuw naar vol.
\end{solution}

\begin{solution}{exo:g7:moon-phases:3}
Nieuwe maan: je kijkt naar de lamp met de bal ernaartoe --- zijn
verlichte helft van je af gekeerd. \omterm{ex:g2:moon-in-the-sky:shapes}{Volle maan}: lamp in je rug, bal
voor je --- zijn verlichte helft helemaal naar je toe gekeerd.
\end{solution}

\begin{solution}{exo:g7:moon-phases:4}
Twee willekeurige: de \omterm{def:g1:light-and-shadows:shadow}{schaduw} van de Aarde wijst van de Zon af en
zou alleen een \omterm{def:g2:moon-in-the-sky:moon}{Maan} in de volle stand kunnen bereiken; een
sikkelmaan staat aan onze zonkant, nergens in de buurt van die
\omterm{def:g1:light-and-shadows:shadow}{schaduw}; en de rand van een ronde \omterm{def:g1:light-and-shadows:shadow}{schaduw} zou nooit de dubbele
kromming van de gibbeuze vorm kunnen uitsnijden.
\end{solution}

\begin{solution}{exo:g7:moon-phases:5}
De nieuwe maan is donker voor ons omdat haar \emph{eigen} nachtkant
naar ons toe staat --- geen \omterm{def:g1:light-and-shadows:shadow}{schaduw} van ons in het spel, en de verre
kant baadt in vol zonlicht. De \omterm{def:g7:shadows-eclipses:lunar}{maansverduistering} is de volle
\omterm{def:g2:moon-in-the-sky:moon}{Maan} die de echte \omterm{prop:g7:shadows-eclipses:umbra}{kernschaduw} van de Aarde kruist: onze
\omterm{def:g1:light-and-shadows:shadow}{schaduw} die werkelijk, en kort, op haar valt.
\end{solution}

\begin{solution}{exo:g7:moon-phases:6}
De volle \omterm{def:g2:moon-in-the-sky:moon}{Maan} staat tegenover de Zon. Terwijl de draaiende Aarde de
Zon onder de westelijke horizon draagt, moet het tegenoverliggende punt
van de hemel --- waar de volle \omterm{def:g2:moon-in-the-sky:moon}{Maan} wacht --- in het oosten opkomen:
\omterm{def:g1:day-and-night:daynight}{zonsondergang} en volle-maansopkomst zijn één gebeurtenis, van twee
kanten bekeken.
\end{solution}

\begin{solution}{exo:g7:moon-phases:7}
De nieuwe maan reist met de Zon mee: om middernacht staan ze allebei
onder je voeten, aan de andere kant van de Aarde. Een nieuwe maan om
middernacht is even onmogelijk als de Zon om middernacht.
\end{solution}

\begin{solution}{exo:g7:moon-phases:8}
Weg van het punt van \omterm{def:g1:day-and-night:daynight}{zonsondergang} --- de boog van de sikkel buigt
naar de Zon toe, de horens ervandaan. De verlichte kant moet naar de
Zon gekeerd zijn, en de horens zijn de wijkende randen van die
verlichte kant: de meetkunde verbiedt ze naar de zon te wijzen.
\end{solution}

\begin{solution}{exo:g7:moon-phases:9}
Tussen eerste kwartier en vol --- ongeveer drie achtste van de weg
rond, aan de avondkant. Ze kwam midden in de middag op. Vijf nachten
later: \omterm{ex:g2:moon-in-the-sky:shapes}{volle maan} --- opkomend bij \omterm{def:g1:day-and-night:daynight}{zonsondergang}, de hele nacht
boven.
\end{solution}

\begin{solution}{exo:g7:moon-phases:10}
\omterm{def:g7:shadows-eclipses:solar}{Zonsverduistering}: alleen bij nieuwe maan --- Zon--Maan--Aarde op
één lijn. \omterm{def:g2:moon-in-the-sky:moon}{Maan}: alleen bij \omterm{ex:g2:moon-in-the-sky:shapes}{volle maan} --- Zon--Aarde--Maan op één
lijn. Zelfs dan zeldzaam, omdat de hellende \omterm{def:g6:sun-earth-moon:axisorbit}{omloopbaan} van de
\omterm{def:g2:moon-in-the-sky:moon}{Maan} haar meestal boven of onder de exacte lijn voert:
\omterm{def:g7:moon-phases:phases}{schijngestalte} is noodzakelijk, uitlijning is niet gegarandeerd.
\end{solution}

\begin{solution}{exo:g7:moon-phases:11}
De Aarde toont \omterm{def:g7:moon-phases:phases}{schijngestalten} tegengesteld aan de onze, op dezelfde
klok: wanneer wij nieuwe maan zien, zien de maanastronauten een
\emph{volle Aarde} --- onze hele daglichtkant naar hen toe gekeerd; bij
onze \omterm{ex:g2:moon-in-the-sky:shapes}{volle maan} zien zij een nieuwe Aarde. Eén meetkunde, gelezen
vanuit de andere stoel.
\end{solution}

\begin{solution}{exo:g7:moon-phases:12}
Pad: zonlicht treft de Aarde, de heldere Aarde verstrooit het
maanwaarts, de nachtkant van de \omterm{def:g2:moon-in-the-sky:moon}{Maan} verstrooit er een beetje van
terug naar onze ogen --- zonlicht dat twee keer is doorgegeven. Bij
onze sikkelschijngestalten ziet de \omterm{def:g2:moon-in-the-sky:moon}{Maan} een bijna
\emph{volle Aarde} (de regel van de tegengestelde \omterm{def:g7:moon-phases:phases}{schijngestalte}): de
nachtkant wordt overgoten door de helderst mogelijke Aarde --- het
sterkste aardschijnsel. Naar onze \omterm{ex:g2:moon-in-the-sky:shapes}{volle maan} toe ziet de
\omterm{def:g2:moon-in-the-sky:moon}{Maan} een dunne sikkelaarde: de schijnwerper slinkt, en de
felheid van de volle \omterm{def:g7:moon-phases:phases}{schijngestalte} verdrinkt wat er over is.
\end{solution}

\begin{solution}{pb:g7:moon-phases:1}
\textbf{1.} Eerste kwartier rond nacht $7$; \omterm{ex:g2:moon-in-the-sky:shapes}{volle maan} rond nacht
$15$; laatste kwartier rond nacht $22$ (de kringloop sluit rond nacht
$29$--$30$).
\textbf{2.} \omterm{ex:g2:moon-in-the-sky:shapes}{Volle maan}: komt op bij \omterm{def:g1:day-and-night:daynight}{zonsondergang}. Eerste
kwartier: komt op rond het middaguur. Laatste kwartier: komt op rond
middernacht. Nieuwe maan: komt op en gaat onder met de Zon.
\textbf{3.} De jonge sikkel staat maar een klein stukje achter de Zon
op de \omterm{def:g6:sun-earth-moon:axisorbit}{omloopbaan}: ze hangt na \omterm{def:g1:day-and-night:daynight}{zonsondergang} in het westen en volgt
de Zon al snel omlaag. De \omterm{ex:g2:moon-in-the-sky:shapes}{volle maan} staat tegenover: als de Zon in
het westen ondergaat, betekent tegenover opkomen in het oosten.
\textbf{4.} Elke \omterm{def:g7:moon-phases:phases}{schijngestalte} die een deel van de dienst van de Zon
deelt: de wassende gibbeuze maan in de loop van de middag, en vooral
de afnemende \omterm{def:g7:moon-phases:phases}{schijngestalten} --- het laatste kwartier en de
afnemende sikkel staan hoog aan de ochtendhemel.
\textbf{5.} \omterm{def:g7:moon-phases:phases}{Wassend} gibbeus: een bal die voor driekwart verlicht is
betekent dat je tussen eerste kwartier en vol staat --- lamp achter je
schouder, bal ervoor en opzij.
\textbf{6.} Ze zien de \emph{eigen} nachtkant van de bal --- de helft
die van de lamp is weggedraaid. De Aarde-jij werpt ook een
\omterm{def:g1:light-and-shadows:shadow}{schaduw}, maar die strekt zich achter je uit, van de lamp af; de bal
zit er nergens in, behalve in de vollemaanstand --- en dan nog alleen
bij exacte uitlijning: een verduistering, geen \omterm{def:g7:moon-phases:phases}{schijngestalte}.
\textbf{7.} De \omterm{def:g2:moon-in-the-sky:moon}{Maan} draait precies één keer per \omterm{def:g6:sun-earth-moon:axisorbit}{omloopbaan}, dus
blijft dezelfde helft voor altijd naar de Aarde gekeerd --- konijn,
gezicht en al; de verre kant wachtte op een camera die eromheen vloog.
\textbf{8.} Draaiing en \omterm{def:g6:sun-earth-moon:axisorbit}{omloopbaan} houden volmaakt gelijke tred:
terwijl de \omterm{def:g2:moon-in-the-sky:moon}{Maan} langs haar pad vordert, draait ze precies genoeg om
hetzelfde gezicht op ons gericht te houden --- als een danser die om
een partner heen cirkelt en hem steeds aankijkt: ze draait echt (de
zaal ziet elke kant), en toont de partner toch maar één gezicht.
\textbf{9.} Opkomen om middernacht is het rooster van het laatste
kwartier --- de afnemende kant: laatste kwartier (de verlichte helft
is naar de richting van de ochtendzon gekeerd).
\textbf{10.} Fouten: een sikkel staat dicht bij de Zon en kan dus
nooit hoog aan de hemel staan om middernacht; horens wijzen \emph{van}
de Zon \emph{af}, nooit ernaartoe; en zonsondergangkleuren om
middernacht zijn een onmogelijkheid op zichzelf. Drie keer mis.
\textbf{11.} Nee: een \omterm{def:g7:shadows-eclipses:solar}{zonsverduistering} vereist dat de \omterm{def:g2:moon-in-the-sky:moon}{Maan}
tussen Zon en Aarde staat --- de nieuwemaanstand. De \omterm{def:g2:moon-in-the-sky:moon}{Maan} van
vanavond staat vol, aan de andere kant: de twee gebeurtenissen kunnen
geen nacht delen.
\textbf{12.} Bijvoorbeeld: ``Eén feit stuurt de maand: de \omterm{def:g2:moon-in-the-sky:moon}{Maan} is
altijd half verlicht, en wij bekijken die helft in een kringloop vanaf
wisselende kanten. Deze maand begraven: de \omterm{def:g1:light-and-shadows:shadow}{schaduw} van de Aarde heeft
niets met de \omterm{def:g7:moon-phases:phases}{schijngestalten} te maken --- die komt alleen bij
zeldzame verduisteringen langs. En neem de klok mee naar huis: elke
\omterm{def:g7:moon-phases:phases}{schijngestalte} houdt haar eigen opkomstuur --- vol bij
\omterm{def:g1:day-and-night:daynight}{zonsondergang}, de kwartieren om het middaguur en om middernacht ---
dus de vorm van de \omterm{def:g2:moon-in-the-sky:moon}{Maan} vertelt je haar rooster.''
\end{solution}
